\chapter{Mathematische Grundlagen}

Dieses Kapitel ist stark an das Skript von
Schaefer-Lorinser~\cite{frank} angelehnt. Hier werden f"ur das
Verst"andnis der Public Key Kryptographie unbedingt notwendigen
Teile aus~\cite{frank} behandelt. F"ur ein besseres Verst"andnis
der Primtests und Faktorisierungsalgorithmen ist die Lekt"ure des
ausf"uhrlichen und weitergehenden Skripts w"armstens empfohlen.

\section{Primzahlen}

In diesem Kapitel werden wir uns mit ein paar grundlegenden
Eigenschaften der Primzahlen besch"aftigen, insbesondere werden wir
ein paar S"atze kennenlernen, die die Verteilung der Primzahlen in den
nat"urlichen Zahlen beschreiben. Beginnen wir jedoch mit der
Definition von Primzahlen:

\begin{definition}[Primzahl]
Jede nat"urliche Zahl $n > 1$, die innerhalb der nat"urlichen Zahlen
nur durch sich selbst und durch 1 teilbar ist, hei\3t {\em Primzahl}.
\end{definition}

Die Primzahlen bilden die fundamentalen Atome der Arithmetik auf den
nat"urlichen Zahlen, da sich jede nat"urliche Zahl in eindeutiger
Weise als Produkt der Primzahlen darstellen l"a\3t. Dies wird auch als
der {\em Hauptsatz der Arithmetik} bezeichnet.

Um auf eine effiziente Weise eine Liste der Primzahlen kleiner einer
vorgegebenen Zahl $n$ zu erstellen, ist das auf den griechischen
Mathematiker Eratosthenes (276-196 v.Chr.) zur"uckgehendes Verfahren
{\em "`Sieb des Eratosthenes"'} geeignet.

Schon aus der Antike ist der von Euklid ca. 325 v.Chr. in seinen
Elementen notierte Satz bekannt:

%%
%% Satz von Euklid
%%
\begin{satz}[von Euklid]
Es gibt unendlich viele Primzahlen.
\end{satz}

W"ahrend uns dieser Satz also zeigt, da\3 in der Folge der
nat"urlichen Zahlen immer wieder eine Primzahl vorkommt, zeigt der
n"achste Satz, da\3 auch beliebig gro\3e L"ucken auftauchen k"onnen,
das hei\3t Abschnitte der nat"urlichen Zahlen, die keine Primzahl
enthalten.

\begin{satz}
Zu jedem $k \in \N$ gibt es in den nat"urlichen Zahlen einen Abschnitt
der L"ange $k$, der keine Primzahl enth"alt.
\end{satz}

\beweis "Ubung \hfill $\Box$

Das hei\3t also, die Primzahlen treten in gewisser Weise
unregelm"a\3ig auf. Es gibt keine geschlossene Formel, die beschreibt,
welche nat"urlichen Zahlen Primzahlen sind. Dennoch kann eine recht
pr"azise Aussage "uber die asymptotische Verteilung der Primzahlen
gemacht werden, wie der folgende von Tschebyscheff bewiesene Satz
zeigt. Wir definieren zun"achst:

\begin{definition}
Es bezeichne f"ur reelle Zahlen $x\in\R$ die Funktion $\pi(x)$ die
Anzahl der Primzahlen, die nicht gr"o\3er als $x$ sind.
\end{definition}
%%
%% Satz von Tschebyscheff
%%
\begin{satz}[Tschebyscheff] 
Es gibt $a,b \in \R_{>0}$, so da\3 f"ur $x\ge 2$ gilt
\begin{equation}
  a\frac{x}{\ln x} < \pi(x) < b\frac{x}{\ln x}.
\end{equation}
\end{satz}

\beweis "Ubung \hfill $\Box$

Betrachten wir einen Restklassenring $\ring{n}$ modulo einer nat"urlichen
Zahl $n$, dann ist unmittelbar klar, da\3 dieser Ring nicht die
K"orperaxiome erf"ullt, falls $n$ keine Primzahl ist. Denn unter
dieser Voraussetzung ist jeder echte Teiler von $n$ nicht
invertierbar. Die teilerfremden Zahlen $a$ mit $1 \le a \le n$, f"ur
die $\ggt(a,n)=1$ gilt, sind in diesem Restklassenring $\ring{n}$
invertierbar. Sie bilden die Einheitengruppe $\einheit{n}$, und die
Anzahl der Elemente dieser Gruppe wird durch die sogenannte Eulersche
$\Phi$-Funktion angegeben.

%%
%% Eulersche Phi-Funktion
%%
\begin{definition}[Eulersche $\Phi$-Funktion] 
  Die Funktion $\Phi(n)$ bezeichnet die Anzahl der nat"urlichen Zahlen 
  kleiner $n$, die relativ prim zu $n$ sind
  und wird die Eulersche $\Phi$-Funktion genannt:
  $$ \#\einheit{n} = \euler{n} = \#\{ m \mid 1 \le m < n, \ggt
(m,n)=1\} $$
\end{definition}   

Wie kann diese Funktion explizit berechnet werden? Betrachten wir
zun"achst den einfachen Fall, da\3 $n$ eine Primzahlpotenz $p^e$ ist.

%%
%% Phi-Funktion f"ur Primpotenz
%%
\begin{bem}[Primzahlpotenz $p^e$]
  Es sei $p$ eine Primzahl und $e \ge 1$. Dann gilt
  $$ \euler{p^e} = p^e - p^{e-1}. $$
\end{bem} 
\beweis
F"ur alle Produkte $s p$ mit $s = 1, \ldots ,p^{e-1}$ gilt $\ggt(s p,
p^e) \ne 1$. Damit sind die $p^e - p^{e-1}$ restlichen Zahlen teilerfremd
zu $p^e$. \hfill $\Box$

Es gilt ferner, da\3 die Eulersche $\Phi$-Funktion f"ur teilerfremde
Zahlen multiplikativ ist:

%%
%% Bem. Multipikativit"at der Phi-Funktion
%%
\begin{bem}[Multiplikativit"at der Eulerschen $\Phi$-Funktion]
  Es seien $m, n \ge 1$ zwei nat"urliche Zahlen mit $\ggt(m,n)=1$. 
  Dann gilt
  $$ \euler{m \cdot n} = \euler{m} \cdot \euler{n} .$$
\end{bem}

\beweis
Da $m, n$ teilerfremd gilt $\ggt(m \cdot n ,s)=1$ genau dann, wenn
$\ggt (m,s)$ und $\ggt(n,s)=1$. Aufgrund der Isomorphie nach dem
Chinesischem Restesatz
$$ \ring{mn} \cong \ring{m} \times \ring{n} $$
ergibt sich die Multiplikativit"at
$$ \euler{m \cdot n} = \euler{m} \cdot \euler{n}$$
unmittelbar, da
$$ \einheit{mn} \cong \einheit{m} \times \einheit{n}. $$
\hfill $\Box$

Aus diesen Bemerkungen folgt unmittelbar das Korollar, das die
$\Phi$-Funktion f"ur alle nat"urlichen Zahlen bestimmt.

%%
%% Kor. Phi-Funktion f"ur beliebige Zahlen
%%
\begin{kor}[Eulersche $\Phi$-Funktion f"ur beliebige Zahlen]
  Jede nat"urliche Zahl $n$ l"a\3t sich eindeutig als $n =
  \prod_{i=1}^{r} p_i^{e_i}$ in seine Primfaktoren zerlegen. Damit
  gilt
  $$ \euler{p_1^{e_1} \cdots p_n^{e_r}} = \prod_{i=1}^r \left(
  p_i^{e_i} - p_i^{e_i-1} \right). $$
\end{kor}

Der folgende Satz geht auf Fermat (1601-1665) zur"uck und wird oft 
als {\em der kleine Satz von Fermat} bezeichnet.

%%
%% Satz kleiner Fermat
%%
\begin{satz}[der kleine Satz von Fermat] \label{fermat}
  Es sei $p$ prim, dann gilt f"ur alle $b\in\einheit{p}$
  $$ b^{p-1} \equiv 1 \pmod{p} .$$
\end{satz}

Das folgende weitergehende Resultat geht auf Euler (1707-1783)
zur"uck. Es stellt eine Verallgemeinerung des kleinen Fermats dar, da
es auch die F"alle, in denen $p$ keine Primzahl ist, erfa\3t.

%%
%% Satz von Euler
%%
\begin{satz} 
Es seien $n,b\in\N$ mit $\ggt(b,n)=1$, das hei"st $b$ ist eine Einheit
in dem Ring $\ring{n}$. Dann gilt
$$ b^{\euler{n}} \equiv 1 \pmod{n} .$$
\end{satz}

Da f"ur $n$ prim gilt: $\euler{n} = n-1$, ist Satz~\ref{fermat} ein
Spezialfall davon.
%%
%% *********** Quadratische Reste
%%

\section{Quadratische Reste}

Wir wollen im folgenden die Frage untersuchen, wann sich eine Einheit $a$
des Restklassenrings $\ring{n}$ als Quadrat einer anderen Einheit $x$
darstellen l"a"st.

\begin{definition}[Quadratische Reste]
  Es seien $a,n \in \N$ mit $\ggt(a,n)=1$. Alle Elemente $a$, f"ur die
  ein $x$ existiert mit
  $$ x^2 \equiv a \pmod{n} $$ hei"sen {\em quadratische Reste} modulo
  $n$, die "ubrigen Einheiten hei"sen {\em quadratische Nicht-Reste}.
\end{definition}

F"ur den Fall, da\3 es sich bei dem Modulus $n$ um eine Primzahl
handelt, definieren wir das {\em Legendre Symbol}:

%%
%% Def. Legendre Symbol
%%
\begin{definition}[Legendre Symbol]
Es sei $p$ eine Primzahl und $a\in\N$. Dann hei\3t das durch 
\begin{equation}
  \leg{a}{p} := \left\{ 
                  \begin{array}{rl}
                        1, & \mbox{falls $a$ ein quadratischer Rest ist} \\
                       -1, & \mbox{falls $a$ kein quadratische Rest ist}
                  \end{array}
                \right.
\end{equation}
definierte Symbol $\leg{a}{p}$ das Legendre Symbol von $a$ nach $p$.
\end{definition}

Der folgende Satz gibt uns schon zwei wichtige Regeln f"ur
quadratische Reste bzw. quadratische Nicht-Reste.

%%
%% Satz
%%
\begin{satz} \label{wannQR}
Das Produkt von zwei quadratischen Resten ist wieder ein quadratischer
Rest, und das Produkt von einem quadratischen Rest und einem
quadratischen Nicht-Rest ist ein quadratischer Nicht-Rest.
\end{satz}

\beweis 
Der erste Teil des Beweises ist direkt einzusehen. Denn f"ur zwei
quadratische Reste $a$ und $b$ exisitieren $x$ und $y$ mit $x^2 \equiv
a \pmod{n}$ und $y^2 \equiv b \pmod{n}$. Daraus folgt aber $(x\cdot
y)^2 \equiv a\cdot b \pmod{n}$.

Den zweiten Fall beweisen wir durch einen Widerspruch. Angenommen, das
Produkt von einem Quadratischen Rest $x^2$ und einem Quadratischen Nicht-Rest 
$z$ w"are ein Quadratischer Rest
$y^2$. Daraus w"urde
$$ x^2 \cdot z \equiv y^2 \pmod{n}$$
folgen. Da $x$ als Einheit invertierbar ist, w"are dann $z$ wegen
$$ z \equiv (x^{-1} \cdot y)^2 \pmod{n}$$
ein Quadratischer Rest, was der Voraussetzung widerspricht. \hfill $\Box$

F"ur den Fall eines Produkts von zwei quadratischen Nicht-Resten gibt
es keine allgemeine Regel. F"ur den Fall aber, da\3 der Modulus $n$
eine Primzahl ist, ergibt dies immer einen quadratischen Rest.

%%
%% Satz Mulitplikativ
%%
\begin{satz} \label{legmultiplikativ}
  Im Falle von quadratischen Resten modulo einer Primzahl $p$ gilt:
  \begin{equation}
    \leg{a}{p} \leg{b}{p} = \leg{ab}{p}. \label{legendreMult}
  \end{equation}
  Ferner ist f"ur $p\ne 2$ die H"alfte der Elemente der Einheitengruppe 
  $\einheit{p}$ Quadratische Reste und die andere H"alfte Quadratische 
  Nicht-Reste.
\end{satz}

\beweis 
F"ur $p=2$ ist die Behauptung offensichtlich. Sei nun $p\ne 2$.  Die
Einheitengruppe $\einheit{p}$ ist eine zyklische Gruppe der Ordnung
$p-1$. Es gibt also einen Erzeuger $g$ der Ordung $p-1$ und
$$ g,g^2,\ldots,g^{p-1} $$ 
ist bereits die ganze $\einheit{p}$. Zu jedem Element
$a\in\einheit{p}$ existiert ein Index $i$ mit
$$ g^i \equiv a \pmod{p}. $$ Der Erzeuger $g$ ist kein Quadratischer
Rest, da sonst ein $x$ mit $x^2 \equiv g \pmod{p}$ existieren
w"urde. Nach dem Satz~\ref{fermat} w"urde f"ur dieses $x$
$$ x^{p-1} \equiv 1 \pmod{p} $$ gelten und somit $g^{(p-1)/2} \equiv 1
\pmod{p}$. Das ist aber ein Widerspruch dazu, da\3 $g$ die ganze
Einheitengruppe $\einheit{p}$ erzeugt.  Offensichtlich sind alle
geraden Potenzen von $g$ Quadratische Reste
$$ g^2,g^4,g^6,\ldots ,g^{p-1}$$
und da $g$ ein Quadratischer Nicht-Rest ist, sind nach Satz~\ref{wannQR}
alle ungeraden Potenzen von $g$ Quadratische Nicht-Reste
$$ g,g^3,g^5, \ldots ,g^{p-2}.$$ 
Damit ist die Formel~\ref{legendreMult} bewiesen, da das
Produkt von zwei ungeraden Potenzen eine gerade Potenz ergibt. 
\hfill $\Box$

Auf Euler geht das folgende Ergebnis zur"uck, das uns ein
Kriterium liefert, f"ur ein gegebenes Element $a\in\einheit{p}$ direkt
zu entscheiden, ob es ein quadratischer Rest oder nicht.
%%
%% ********* Euler Kriterium
%%
\begin{satz}[Euler Kriterium] 
Es sei $a\in\N$ mit $\ggt(a,p)=1$, wobei $p$ eine ungerade Primzahl
ist. Dann gilt
$$ \leg{a}{p} \equiv a^{\frac{p-1}{2}} \pmod{p}. \label{eulerkrit} $$
\end{satz}

\beweis 
Es sei $g$ ein Erzeuger der Einheitengruppe $\einheit{p}$. Wir
bestimmen zun"achst ein $s\in\left\{0,\ldots,p-2\right\}$ mit 
$$ g^s \equiv -1 \pmod{p} $$. 
F"ur ein solches $s$ gilt 
$$ g^{2s} \equiv (g^s)^2 \equiv (-1)^2 \equiv 1 \pmod{p}. $$ 
Dies kann nur gelten, wenn $2s$ ein Vielfaches
der Gruppenordnung $\#\einheit{p} = p-1$ ist, da $p-1$ der kleinste
Exponent $e$ ist, f"ur den $g^e \equiv 1 \pmod{p}$ gilt. Also kann $s$
nur einen der Werte
$$ 0,{p-1 \over 2}, p-1,\ldots $$
annehmen. Es gilt wegen Satz~\ref{fermat}
$$ g^0 \equiv g^{p-1} \equiv 1 \pmod{p}. $$
Also mu\3 $s={p-1 \over 2}$ sein und somit 
$$ g^{(p-1)/2} \equiv -1 \pmod{p}. $$ 
Da $g$ ein Erzeuger der Einheitengruppe $\einheit{p}$ ist, kann jedes Element
$a\in\einheit{p}$ als $a \equiv g^i \pmod{p}$ dargestellt
werden. Damit gilt
$$ a^{(p-1)/2} \equiv (g^i)^{(p-1)/2} \equiv
\left(g^{(p-1)/2}\right)^i \equiv (-1)^i \pmod{p}. $$ Das hei\3t aber,
da\3 $a^{(p-1)/2} \equiv 1 \pmod{p}$ genau dann gilt, wenn $i$ ein
gerader Exponent ist, und $a^{(p-1)/2} \equiv -1 \pmod{p}$ genau dann,
wenn $i$ ein ungerader Exponent ist. Dies ist aber das Kriterium f"ur Quadratische
Reste aus Satz~\ref{legmultiplikativ}. \hfill $\Box$

\begin{kor} \label{modulo}
Es sei $p$ eine ungerade Primzahl und es gelte
$\ggt(a,p)=\ggt(b,p)=1.$ Aus $a \equiv b \pmod{p}$ folgt, da\3
$\leg{a}{p} = \leg{b}{p}$.
\end{kor}

\beweis 
Es ist $a \equiv b \pmod{p}$ gleichbedeutend damit, da\3 ein $l\in \Z$
existiert mit $b=a+l\cdot p$. Da modulo $p$ gerechnet wird, gilt:
\begin{eqnarray*}
  \leg{b}{p} & \equiv & \leg{a+l\cdot p}{p} \\
             & \equiv & (a+l\cdot p)^{(p-1)/2} \\
             & \equiv & \sum_{i=0}^{p-1} {p-1 \choose 2i} a^{(p-1)/2-i}(lp)^i \\
             & \equiv & a^{(p-1)/2} \pmod{p}.
\end{eqnarray*}
\rule{0pt}{0pt}\hfill $\Box$
                     
Ein sehr interessantes Ergebnis der Zahlentheorie, das wir ohne Beweis
angeben, ist {\em das quadratische Reziprozit"atsgesetz}, das von
Gau\3 bewiesen wurde. Es stellt eine Relation zwischen dem
Legendresymbol von zwei Primzahlen $\leg{p}{q}$ und der "`Umkehrung"'
$\leg{q}{p}$ dar.

%%
%% Quadratisches Reziprozit"atsgesetz
%%
\begin{satz}[Quadratisches Reziprozit"atsgesetz] 
Seien $p$ und $q$ zwei ungerade Primzahlen. Dann gilt:
  \begin{equation}
    \leg{p}{q} = \leg{q}{p} (-1)^{{1 \over 2}(p-1){1 \over 2}(q-1)} .
  \end{equation}
\end{satz}

Mit Hilfe des quadratischen Reziprozit"atssatzes und des
Korollars~\ref{modulo} l"a\3t sich effizient entscheiden, ob eine
gegebene Zahl $a$ modulo einer Primzahl $p$ ein Quadratischer Rest ist.

%%
%% Beispiel
%%
\begin{beispiel} Ist 19 ein Quadratischer Rest modulo 61? Es gilt
$$\leg{19}{61}=\leg{61}{19}(-1)^{30 \cdot 9}=\leg{61}{19}=\leg{4}{19} = 1, $$
da $4 \equiv 2^2 \pmod{19}$.
\end{beispiel}

%%
%% ***************** Jacobi Symbol
%%
Beim Legendre-Symbol sind Quadratische Reste modulo einer Primzahl $p$
zugelassen. Im folgenden wird eine Verallgemeinerung des
Legendre-Symbols, das Jacobi-Symbol, das auch f"ur zusammengesetzte
Zahlen definiert ist, eingef"uhrt.

%%
%% Def.
%%
\begin{definition}[Jacobi-Symbol]
  Es seien $a,n\in\N$ mit $n$ ungerade und $\ggt(a,n)=1$. Die
  Primfaktorzerlegung von $n$ sei gegeben durch $n=\prod_{i=1}^{r}
  p_i^{e_i}$. Dann ist das Jacobi-Symbol definiert durch
  $$\leg{a}{n} := \prod_{i=1}^{r} \leg{a}{p_i}^{e_i} .$$
\end{definition}

Das Jacobi-Symbol ist so definiert, da\3 es den gleichen Rechenregeln
wie das Legendre-Symbol gen"ugt. Es ist aber zu beachten, da\3 das
Jacobi-Symbol $\leg{a}{n}$ {\em nicht} angibt, ob $a$ ein
quadratischer Rest modulo $n$ ist.

Es ist leicht einzusehen, da\3 f"ur ein $a$, das ein Quadratischer
Rest modulo einer zusammengesetzten Zahl $n=\prod_{i=1}^{r}p_i^{e_i}$
ist, das Jacobi-Sybmol $\leg{a}{n}$ den Wert 1 annimmt. Da $a$ ein
Quadratischer Rest modulo $n$ ist, existiert ein $x\in\N$ mit $x^2
\equiv a \pmod{n}$. Betrachten wir diese Gleichung modulo $p_i$ f"ur
$i=1,\ldots,r$, so erhalten wir $ x^2 \equiv a \pmod{p_i}$. Es ist
also $a$ auch Quadratischer Rest modulo aller Primfaktoren $p_i$ 
f"ur $i=1,\ldots,r$ und damit ist das Jacobi-Symbol $\leg{a}{n}$ 
gleich 1.

Die Umkehrung dieser Aussage gilt nicht, wie das folgende Beispiel zeigt.

\begin{beispiel}
  Das Jacobi-Symbol kann auch f"ur Quadratische Nicht-Reste den Wert 1 annehmen:
  $$ \leg{2}{15} = \leg{2}{3}\leg{2}{5} = (-1)(-1)=1, $$
  obwohl 2 kein Quadratischer Rest modulo 15 ist.
\end{beispiel}    

%%
%% ************** Primtest
%%
\section{Primtests}

Die Primzahltests weisen eine ganze Reihe von Methoden und Parallelen
zu den Faktorisierungsalgorithmen auf

\subsection{Pseudoprimzahlen}

Betrachten wir zun"achst einen einfachen Satz, der ein Kriterium
liefert, mit dem sich entscheiden l"a\3t, ob eine Zahl $n$ prim oder
zusammengesetzt ist.

%%
%% Satz Prim-Kriterium
%%
\begin{satz}
   Eine Zahl $n$ ist prim genau dann, wenn
   \begin{equation} 
     a^{n-1} \equiv 1 \pmod{n}, \label{primkrit} 
   \end{equation}
   f"ur alle $a\in\N$ mit $2 \le a < n$.
\end{satz}

\beweis 
Die Hinrichtung gilt aufgrund des Satzes~\ref{fermat}.

F"ur die R"uckrichtung nehmen wir an, die
Fermat-Bedingung~\ref{primkrit} sei erf"ullt und $n$ sei
zusammengesetzt. Dann hat $n$ einen nichttrivialen Teiler $r$ mit $1 <
r < n$ und $n=r\cdot s$, und es gilt
$$ r^{n-1} \equiv 1 \pmod{n}. $$
Es existiert also ein $l\in\Z$ mit
$$ 1=r^{n-1}+l\cdot n = r^{n-1}+l\cdot s\cdot r = r^{n-1} + (l\cdot s) \cdot r. $$
Betrachten wir diese Gleichung modulo $r$, so ergibt sich 
$$ 0 \equiv 1 \pmod{r}. $$ 
Das ist ein Widerspruch. \hfill $\Box$

Mit Hilfe dieses Satzes k"onnen wir bereits einen einfachen Primtest realisieren.
Haben wir eine Zahl $n$ gegeben, von der wir entscheiden m"ochte, ob sie prim
oder zusammengesetzt ist, so w"ahlen wir eine Basis $a$ mit
$2\le a < n$ und testen die Fermat-Bedingung~\ref{primkrit}.
Erf"ullt die Zahl $n$ die Bedingung nicht, so ist $n$ sicher
eine zusammengesetzte Zahl. Andernfalls w"ahlen wir eine neue Basis
und f"uhren den Test erneut durch.

Dies w"urde, falls $n$ eine Primzahl ist, einen extrem aufwendigen
Primtest ergeben. Wir k"onnen aber nach einer gr"o\3eren Anzahl von
Durchl"aufen des Primtests schon mit einer gewissen Wahrscheinlichkeit
schlie\3en, da\3 die Zahl $n$ prim ist.

\begin{definition}
  Eine zusammengesetzte Zahl $n$ hei\3t {\em Pseudoprimzahl zur Basis $a$}, falls
  $$ a^{n-1} \equiv 1 \pmod{n} $$
  erf"ullt ist.
\end{definition}

\begin{beispiel}  
  Die Zahl $341=11\cdot 31$ ist eine Pseudoprimzahl zur Basis 2, denn
  $$ 2^{341-1} \equiv 1 \pmod{341}. $$
\end{beispiel}

\begin{satz}
  Es sei $n\in\N$ eine Pseudoprimzahl zur Basis $a_1$ und zur Basis
  $a_2$ mit $a_1, a_2\in\einheit{n}$. Dann ist $n$ auch eine
  Pseudoprimzahl zur Basis $a_1 \cdot a_2$ und zur Basis $a_1 \cdot
  a_2^{-1}$.

  Falls ein $a\in \einheit{n}$ existiert, das die Bedingung
  $$ a^{n-1} \equiv 1 \pmod{n} $$ nicht erf"ullt, so wird sie auch von
  mindestens der H"alfte der Elemente aus $\einheit{n}$ nicht
  erf"ullt.
\end{satz}

\beweis 
Der erste Teil ist durch direktes Nachrechnen zu beweisen. Es gilt
$$(a_1\cdot a_2)^{n-1} \equiv a_1^{n-1} \cdot a_2^{n-1} \equiv 1
\pmod{n},$$ und damit ist $n$ eine Pseudoprimzahl zur Basis $a_1
a_2$. Weiterhin gilt
$$(a_1 \cdot a_2^{-1}) \equiv a_1^{n-1} \cdot (a_2^{n-1})^{-1} \equiv
1 \cdot 1^{-1} \equiv 1 \pmod{n},$$ und damit ist $n$ eine
Pseudoprimzahl zur Basis $a_1 a_2^{-1}$.

Es sei $B:=\{a_1,\ldots, a_l\}$ die Menge aller $a\in\einheit{n}$ mit
$a^{n-1} \equiv 1 \pmod{n}$.  Die zusammengesetzte Zahl $n$ ist also
eine Pseudoprimzahl zur Basis $a_i\in B$ f"ur $i=1,\ldots, l$. Wir
w"ahlen $b\in\einheit{n}\setminus B$ beliebig aus. Wenn $n$ eine
Pseudoprimzahl zur Basis $b \cdot a_i$ w"are, dann w"are nach dem
ersten Teil $n$ auch eine Pseudoprimzahl zur Basis $(b \cdot a_i)
a_i^{-1} \equiv b \pmod{n}$. Dies kann aufgrund der Wahl von $b$ aber
nicht sein, und damit sind alle Elemente aus $\{ b \cdot a_1, \ldots,
b \cdot a_n\}$ Basen, zu denen $n$ keine Pseudoprimzahl ist. Damit
sieht man, da\3 mindestens die H"alfte der Elemente aus $\einheit{n}$
Basen sind, zu denen $n$ keine Pseudoprimzahl ist. \hfill $\Box$

Falls also zu einer zusammengesetzten Zahl $n$ "uberhaupt eine Zahl
$a\in\einheit{n}$ existiert, so da\3 $n$ keine Pseudoprimzahl zur
Basis $a$ ist, so finden wir mit einer Wahrscheinlichkeit gr"o\3er als
$1/2$ eine solche Basis. Eine solche Basis $a$ wird auch {\em Zeuge}
f"ur $n$ genannt, weil die Angabe von $a$ gen"ugt, um zu beweisen,
da\3 $n$ eine zusammengesetzte Zahl ist.

W"ahlen wir nun $k$ verschiedene Kandidaten f"ur die Basis $a$ auf
zuf"allige Weise, so ist $a$ f"ur eine zusammengesetzte Zahl $n$ nur mit
Wahrscheinlichkeit kleiner als $1/2^{k}$ unter diesen Elementen kein
Zeuge. Diese Aussage gilt wieder nur unter der Voraussetzung, da\3 es
f"ur $n$ "uberhaupt einen Zeugen gibt.

Es gibt Zahlen, f"ur die das nicht erf"ullt ist. Sie werden {\em
Carmichael-Zahlen} genannt. Wenn $n$ eine Carmichael-Zahl ist, dann
gilt
$$ a^{n-1} \equiv 1 \pmod{n}$$ f"ur alle $a\in\einheit{n}$.

Die Zahl $561=3\cdot 11\cdot 17$ ist beispielsweise eine Carmichael-Zahl.

\begin{satz}
  Es seien $m,n\in\Z$ ungerade. F"ur das Jacobi-Symbol gelten folgende
  Rechenregeln: \\
  (i)   $$ \leg{-1}{m} = (-1)^{(m-1)/2}, $$
  (ii)  $$ \leg{2}{m} = (-1)^{(m^2-1)/8}, $$
  (iii) F"ur $m,n$ teilerfremd gilt:
  $$ \leg{m}{n} = \leg{n}{m} (-1)^{\frac{m-1}{2}\frac{n-1}{2}}. $$
\end{satz}


%% ******** Primzahltests
\subsection{Probabilistischer Primtest nach Solovay, Strassen}

Im vorigen Abschnitt haben wir aus dem kleinen Fermat eine Art
heuristischen Primtest abgeleitet, indem wir f"ur die auf Primalit"at
zu testende Zahl $n$ "uberpr"uften, ob sie die Eigenschaft~\ref{primkrit}, die der
kleine Fermat f"ur alle Primzahlen liefert, erf"ullt. Dieser einfache
Primtest versagt mit sehr hoher Wahrscheinlichkeit bei den
Carmichael-Zahlen. Wir w"ahlen in dem Primtest eine Zufallszahl $a$ mit
$1 < a < n$. Ist $\ggt(a,n)\neq 1$, so hat man einen Nullteiler
gefunden, und wei\3 sofort, da\3 die Zahl $n$ zusammengesetzt ist. Wie
man aber leicht aus der Berechnungsvorschrift f"ur die Eulersche
$\Phi$-Funktion sieht, ist die Menge der Nullteiler im Vergleich zu
der Menge der Einheiten sehr klein. Ist $n$ eine Carmichael-Zahl, so
ist $n$ eine Pseudoprimzahl zu allen Einheiten aus $einheit{n}$.

Die Grundidee des folgenden Primtests ist es, einen analogen Test
bez"uglich des Euler-Kriteriums vorzunehmen.

%%
%% Def.
%%
\begin{definition}
  Analog zur Definition von Pseudoprimzahlen nennen wir
  zusammengesetzte Zahlen $n$, f"ur die
  $$ \leg{a}{n} \equiv a^{(n-1)/2} \pmod{n} $$ gilt, {\em
  Euler-Pseudoprimzahlen} zur Basis $a$.
\end{definition}

Das folgende einfache Lemma zeigt schon, da\3 dieser Primtest den auf
dem kleinen Fermat basierenden Primtest umfa\3t.

%%
%% Lemma
%%
\begin{lemma}
  Wenn $n$ eine Euler-Pseudoprimzahl zur Basis $a$ ist, dann ist $n$
  auch eine Pseudoprimzahl zur Basis $a$.
\end{lemma}

\beweis 
Quadrieren der Gleichung im Satz~\ref{eulerkrit} liefert die
Behauptung. \hfill $\Box$

Der auf dieser "Uberlegung basierende Primtest wurde 1974 von Solovay
und Strassen vorgeschlagen. Man nennt Algorithmen dieser Art
probabilistisch, da es zusammengesetzte Zahlen gibt, f"ur die der Test
behauptet, die Zahl $n$ sei mit einer gewissen Wahrscheinlichkeit
prim. Durch Wiederholen des Primtests kann die
Irrtumswahrscheinlichkeit exponentiell verringert werden. In der
Praxis treten Zahlen, f"ur die sich der Test auch nach mehrmaligem
Anwenden irrt, "au\3erst selten auf.

Der Primzahltest nach Solovay und Strassen war das erste
probabilistische Verfahren "uberhaupt. Heute werden in vielen
Bereichen der Algorithmentechnik probabilistische Methoden eingesetzt.

\subsection*{Algorithmus:}
  \begin{enumerate}
    \item Bestimme eine zuf"allig gew"ahlte Zahl $a$ mit $1 < a < n$.
    \item Teste, ob $\leg{a}{n} \equiv a^{(n-1)/2} \pmod{n}$ gilt.
    \item Falls ja, dann ist $n$ wahrscheinlich prim, sonst ist $n$
          mit Sicherheit zusammengesetzt.
  \end{enumerate}

Das Jacobi-Symbol kann hierbei offensichtlich nicht "uber die
Definiton berechnet werden, da dies die Primfaktorzerlegung von $n$
voraussetzt. Es l"a"st sich effizient mit den Rechenregeln ausrechnen.

Der gro\3e Vorteil dieses Primtests liegt nun darin, da\3 es f"ur
Euler-Pseudoprimzahlen nichts "Aquivalentes zu den Carmichael-Zahlen
gibt. Ferner l"a\3t sich zeigen, da\3 die Irrtumswahrscheinlichkeit
dieses Test kleiner als $1/2$ ist.

\begin{satz}[Solovay, Strassen 1974] 
Es sei $P:=\left\{a \mid \ a\in\einheit{n},\leg{a}{n}\equiv
a^{(n-1)/2}\right\}$.
  \begin{itemize}
    \item Falls $n$ prim ist, so gilt $\#P=n-1$.
    \item Falls $n$ nicht prim ist, so gilt $\#P \le {1 \over 2}(n-1)$.
  \end{itemize}
\end{satz}
\beweis 
Falls $n$ eine Primzahl ist, so liefert das Euler-Kriterium die
Richtigkeit der Behauptung. Es sei nun $n$ eine zusammengesetzte Zahl
mit der Primfaktorzerlegung $n=\prod_{i=1}^{r} p_i^{e_i}$. Wir
definieren zur Vereinfachung zwei Abbildungen

\begin{equation}
  \epsilon : \left\{ \begin{array}{ccc}
                      \einheit{n} & \rightarrow & \einheit{n} \\
                      a           & \mapsto         & \leg{a}{n}
                    \end{array} \right.
\end{equation}
und 
\begin{equation}
  \delta   : \left\{ \begin{array}{ccc}
                      \einheit{n} & \rightarrow & \einheit{n} \\
                      a           & \mapsto         & a^{(n-1)/2}
                    \end{array} \right. .
\end{equation}

Es ist also zu zeigen, da\3 $\#\{a\in\einheit{n} \mid \epsilon(a) =
\delta(a)\} \le {1 \over 2}(n-1)$ erf"ullt ist. Wir machen hierzu eine
Fallunterscheidung.

\begin{enumerate}
%% 1. Fall
\item Es gibt ein $j$ mit $e_j > 1$ ($n$ ist nicht quadratfrei).

      Angenommen es sei $\epsilon(a)=\delta(a)$ f"ur alle 
      $a\in\einheit{n}$. Dann gilt 
      \begin{equation}
        a^{n-1} \equiv \left(a^{(n-1)/2}\right)^2 \equiv \delta(a)^2
        \equiv \epsilon(a)^2 \equiv \leg{a}{n}^2 \equiv 1 \pmod{n}
        \label{gleich}
      \end{equation}
      f"ur alle $a\in\einheit{n}$. Es gilt ferner 
      \begin{eqnarray*}
        \#\einheit{n} & = & \euler{n} \\
                      & = & \euler{p_1^{e_1}} \cdots \euler{p_r^{e_r}} \\
                      & = & p_1^{(e_1-1)}(p_1-1) \cdots p_r^{(e_r-1)}(p_r-1).
      \end{eqnarray*}

      Daraus folgt $p_j \mid \#\einheit{n}$, da $e_j > 1$. Also gibt
      es ein Element $a\in\einheit{n}$ der Ordnung $p_j$, d.h. $p_j$
      ist der kleinste Exponent, f"ur den $a^{p_j} \equiv 1 \pmod{n}$
      gilt. Da $p_j$ kein Teiler von $n-1$ ist, gilt $a^{n-1}
      \not\equiv 1 \pmod{n}$.  Dies ist aber ein Widerspruch zu
      Gleichung~\ref{gleich}. Es ist also $\epsilon \neq \delta$ und
      damit ist die Menge
      $\{a\in\einheit{n}\mid\epsilon(a)=\delta(a)\}$ eine echte
      Untergruppe (sie ist multiplikativ und enth"alt das neutrale
      Element) von $\einheit{n}$.  Damit gilt f"ur die
      Untergruppenordung
      $$\#\{a\in\einheit{n}\mid\epsilon(a)=\delta(a)\} \le {1 \over 2}\euler{n} $$
      und somit die Behauptung, da $\euler{n} \le n-1$ immer erf"ullt ist.
%% 2. Fall
\item Es gilt $e_i=1$ f"ur alle $i=1,\ldots, r$ ($\ring{n}$ ist isomorph 
      zu dem direkten Produkt der Primk"orper $\ring{p_1},\ldots ,\ring{p_r}$).

      Angenommen, es sei $\epsilon(a)=\delta(a)$ f"ur alle
      $a\in\einheit{n}$. W"ahle $a_1$ mit $\leg{a_1}{p_1}=-1$ und
      $a_i$ mit $\leg{a_i}{p_i}=1$ f"ur
      $i=2,\ldots,r$. Nach dem Chinesischem Restesatz existiert ein
      $a\in\Z$ mit $a\equiv a_i \pmod{p_i}$ f"ur $i=1,\ldots,r$. Es
      gilt nach Definition des Jacobi-Symbols und
      Korollar~\ref{modulo}
      $$ \epsilon(a) = \leg{a}{n} = \prod_{i=1}^{r} \leg{a}{p_i} = 
         \prod_{i=1}^r \leg{a_i}{p_i} = -1 \cdot 1^{r-1} = -1 $$
      und nach der Annahme ist $\delta(a)=-1$. Betrachten wir den 
      durch den Chinesischen Restesatz gegebenen Isomorphismus
      $$ \phi :  \ring{n} \longrightarrow \ring{p_1} \times \ldots 
         \times \ring{p_r}.$$

      Da $\phi$ insbesondere ein Homomorphismus ist und aus $x \equiv
      -1 \pmod{n}$ f"ur alle Primteiler $p_i$ von die Gleichungen $x
      \equiv -1 \pmod{p_i}$ folgen, gilt speziell

      \begin{equation} 
         -1 = \delta(a) = a^{(n-1)/2} \mapsto 
         \left(a_1^{(n-1)/2},\ldots,a_r^{(n-1)/2}\right) = (-1,\ldots,-1). 
         \label{minuseins} 
      \end{equation}

      Nach dem Chinesischen Restesatz existiert ein $b\in\Z$ mit
      $b\equiv 1 \pmod{p_1}$ und $b\equiv a_i \pmod{p_i}$ f"ur
      $i=2,\ldots,r$. Damit gilt
      $$ \delta(b) = b^{(n-1)/2} \mapsto \left( 1^{(n-1)/2},
         a_2^{(n-1)/2},\ldots,a_r^{(n-1)/2)} \right) =
         (1,-1,\ldots, -1). $$

      Es kann also $b^{(n-1)/2} \equiv \pm 1$ auf keinen Fall gelten,
      da $b^{(n-1)/2}$ durch den Isomorphismus $\phi$ auf
      $(1,-1,\ldots,-1)$ abgebildet wird und da aufgrund der
      Eigenschaft des neutralen Elements und der
      Formel~\ref{minuseins} gilt

      \begin{eqnarray*}
         1 & \mapsto & \phi(1)  = (1,1,\ldots,1) \\
        -1 & \mapsto & \phi(-1) = (-1,-1,\ldots,-1).
      \end{eqnarray*}

      Das Jacobi-Symbol $\leg{b}{n}$ kann nur die Werte $1$ oder $-1$ annehmen. 
      Es gilt aber $\leg{b}{n}=\epsilon(b)\neq \pm 1$. Also ist auch 
      $\epsilon \neq \delta$ in diesem Fall gezeigt. Mit
      dem gleichen Untergruppenkriterium wie in ersten Fall folgt die
      Behauptung. \hfill $\Box$
\end{enumerate}


%% ************* Rabin und Miller
\subsection{Probabilistischer Primtest nach Rabin und Miller}

Der folgende Primtest wurde von Rabin aufbauend auf den
Ideen von Miller entwickelt. Er liefert eine Verbesserung des
Primtests von Solovay und Strassen, da die Menge der F"alle, in
denen er sich irrt, eine echte Untermenge der F"alle ist, in denen sich
der Primtest von Solovay und Strassen irrt.

Der Algorithmus geht von folgender grunds"atzlicher "Uberlegung aus:
Nehmen wir an, wir h"atten eine Pseudoprimzahl $n$ zur Basis $a$, also
mit
$$ a^{n-1} \equiv 1 \pmod{n}, $$ 
gegeben. Betrachten wir Quadratwurzeln aus $a^{n-1}$, so w"urde f"ur eine 
Primzahl $n$ gelten, da\3 die Quadratwurzel $\pm 1$ sind, da in einem K"orper 
zwei L"osungen f"ur $X^2 \equiv 1 \pmod{n}$ existieren. Das hei\3t, f"ur 
eine Primzahl $n$ w"urde die Folge
$$ a^{(n-1)/2}, a^{(n-1)/2^2},\ldots,a^{(n-1)/2^k} $$ 
f"ur $n-1 = 2^k m$ mit $m$ ungerade aus entweder beliebig vielen 1en, gefolgt von
einer $-1$ und dann beliebigen Zahlen aus $\einheit{n}$ oder lauter
1en bestehen. Diese Beobachtung wird in dem Miller-Rabin-Test
ausgenutzt. Zur Berechnung geht man hier nat"urlich von
$a^{(n-1)/2^k}$ aus und quadriert.

\subsection*{Algorithmus}

\begin{enumerate}
  \item Schreibe $n-1 = 2^k \cdot m$ mit $m$ ungerade.
  \item W"ahle eine Zufallszahl $a$ mit $1<a<n$ und $\ggt(a,n)=1$.
  \item Initialisiere $i:=0$ und $b:=a^m$.
  \item Falls $(i=0$ und $b=1)$, dann ist $n$ wahrscheinlich prim. Stop
  \item Falls $b=-1$, dann ist $n$ wahrscheinlich prim. Stop
  \item Falls $(i>0$ und $b=1)$, dann ist $n$ sicher zusammengesetzt. Stop \\
        Sonst $b:=b^2$ und $i:=i+1$
  \item Falls $i=k+1$, dann ist $n$ sicher zusammengesetzt. Stop \\ 
        Sonst gehe zu Schritt 5.
\end{enumerate}

Der folgende Satz liefert eine Absch"atzung der
Fehlerwahrscheinlichkeit dieses Algorithmus.

%%
%% Satz Rabin
%%
\begin{satz}[Rabin]
  Es sei $n\in\N$ ungerade und mit $n-1=2^k \cdot m$, wobei $m$
  ungerade ist. Zu $a\in\N$ mit $0<a<n$ definieren wir eine Folge
  $a_0,\ldots,a_k$ mit
  $a_i := a^{2^i \cdot m}$ f"ur $0\le i\le k$.
  Es sei 
    $$ P:=\left\{a\in\einheit{n} \mid (a_0 \equiv 1 \pmod{n}) \mbox{
       oder } (\exists i<k : a_i \equiv -1 \pmod{n}\right\}. $$
  \begin{itemize}
    \item Falls $n$ prim, so gilt $\#P = n-1$.
    \item Falls $n$ nicht prim, so gilt $\#P \le \frac{1}{4}(n-1)$.
  \end{itemize}
\end{satz}  

\beweis 
siehe \cite{frank} \hfill $\Box$

Der Algorithmus erkennt die Primzahlen mit einer Wahrscheinlichkeit
gr"o\3er als 3/4. Durch $t$-faches Wiederholen des Algorithmus kann
die Fehlerwahrscheinlichkeit auf $1/4^t$ gesenkt werden.

Der n"achste Satz zeigt, da\3 der Algorithmus von Rabin auch wirklich
eine Verbesserung des Algorithmus von Solovay und Strassen darstellt.
\begin{satz} 
  Es sei 
  $$ P':=\left\{a\in\einheit{n} \mid \leg{a}{n}\equiv a^{(n-1)/2}\right\}$$
  und
  $$ P:=\left\{a\in\einheit{n} \mid (a_0 \equiv 1 \pmod{n}) \mbox{ oder } 
     (\exists i<k : a_i \equiv -1 \pmod{n}\right\}, $$
  wobei $0<a<n$, $n$ ungerade und zusammengesetzt mit $n-1 = 2^k \cdot m$ und 
  $(a_i)_i$ eine 
  Folge mit $a_0,\ldots, a_k$. Dann gilt f"ur zusammengesetzte
  Zahlen $n$: $P \subset P'$. 
\end{satz}

\beweis 
siehe \cite{frank} \hfill $\Box$

\section{Faktorisierungsalgorithmen mit Gruppen}

Den beiden in diesem Abschnitt vorgestellten Algorithmen liegt eine
gemeinsame Idee zugrunde. Der chinesische Restesatz besagt, da\3 f"ur
zusammengesetztes $n$ der Restklassenring $\ring{n}$ isomorph ist zu einem 
direkten Produkt von kleineren Restklassenringen.  Wir definieren eine 
algebraische Struktur, die von dem Restklassenring $\ring{n}$ abh"angt. Wenn wir
in dieser Struktur rechnen, so rechnen wir implizit auch in der algebraischen
Struktur, die gewisserma\3en von den kleineren Komponenten des
direkten Produkts abh"angt, und erh"alten aufgrund des kleinen Fermats
Information "uber eine der kleinere Strukturen und k"onnen dadurch einen
Faktor von $n$ abspalten. Im Falle der $p-1$-Methode ist die
algebraische Struktur $\ring{n}$ selbst und bei der Faktorisierung mit
elliptischen Kurven eine elliptische Kurve $E(\ring{n})$ "uber dem
Restklassenring $\ring{n}$.

\subsection{Pollard $p-1$-Methode}

Der Aufwand dieses Algorithmus h"angt im wesentlichen von der Gr"o\3e
des zu findenden Faktors ab. Nehmen wir also an, die Zahl $n$, die
faktorisiert werden soll, enth"alt einen Primfaktor $p$, der die
Eigenschaft besitzt, da\3 $p-1$ in kleine Faktoren zerf"allt. Konkret
sollen alle Faktoren von $p-1$ kleiner als eine Schranke $k$ sein, und
es soll dar"uberhinaus gelten, da\3 $k!$ durch $(p-1)$ teilbar ist.

Da die Exponentiation modulo $n$ effizient berechnet werden kann, kann
die Berechnung von
$$ m:\equiv 2^{k!} \pmod{n} $$ f"ur nicht zu gro\3e $k$ und $n$
durchgef"uhrt werden. Der kleine Fermat liefert wegen $p-1 \mid k!$,
da\3
$$ m \equiv 2^{k!} \equiv 2^{k! \pmod{p-1}} \equiv 2^0 \equiv 1
\pmod{p} $$ gilt. Damit teilt $p$ also $m-1$. Falls $k$ nicht zu gro\3
angesetzt ist, teilt $n$ mit einer hohen Wahrscheinlichkeit $m-1$
nicht, so da\3 der gr"o\3te gemeinsame Teiler
$$ g:=\ggt(m-1,n) $$
den nicht trivialen Teiler $p$ von $n$ liefert.

Bei der praktischen Anwendung dieser Vorgehensweise, wei\3 man
nat"urlich nicht, in welcher Gr"o\3enordung der richtige Parameter $k$
gew"ahlt werden sollte. Damit nicht alle Primfaktoren von $n$ erfa\3t
werden, ist es wichtig von Zeit zu Zeit $\ggt(m-1,n)$ zu
berechnen. Falls der ggT gleich $1$ ist, so erh"oht man $k$ und
rechnet weiter.

Sollte der ggT schon $n$ sein, so sind alle Primfaktoren von $n$ schon
erreicht worden, so da\3 man entweder mit kleinerem $k$ oder mit einer
anderen Basis als 2 einen neuen Versuch macht.

Der Algorithmus f"uhrt nicht zum Erfolg, falls die Primfaktoren $p_i$ 
von $n$ eine "ahnliche Zerlegung f"ur $p_i-1$ haben. Dann gilt
$(p_i-1)\mid k!$. In diesem Fall wird ganz $n$ abgespalten.

\subsubsection{Verfeinerung des Algorithmus}

Anstatt $k!$ als Exponent zu verweden, ist es g"unstiger eine etwas
andere Wahl zu treffen, denn f"ur sehr kleine Primzahlen $q$ taucht
$q$ in $k!$ mit einem sehr hohem Exponenten auf. Es ist aber
unwahrscheinlich, da\3 dies f"ur die Bedingung $p-1\mid k!$ ben"otigt
wird. In der Praxis kann man beispielsweise so vorgehen, da\3 man bis
zu einer gewissen Schranke $B_1$, alle Primzahlen und Primpotenzen
kleiner $B_1$ betrachtet. An Stelle von $k!$ w"ahlt man dann das Produkt
dieser Primzahlen und Primpotenzen. Das Produkt sei $l$.

Man w"ahlt eine weitere Schranke $B_2 > B_1$ und testet
zus"atzlich zu 
$$ \ggt(2^l-1,n) $$
auch noch
$$ \ggt(2^{l \cdot q_i}-1,n) $$
f"ur alle Primzahlen $q_i$ mit $B_1 < q_i < B_2$, um
die Primfaktoren zu trennen.

\subsubsection{Laufzeitbetrachtungen}

Die Laufzeit dieses Algorithmus h"angt nicht nur von der Gr"o\3e der
zu faktorisierenden Zahl $n$ ab, sondern wird auch stark von der
Zerlegung von $p-1$ f"ur alle Primfaktoren $p$ beeinflu\3t. Die
Primfaktorzerlegung von $n$ ist im voraus nicht bekannt. Daher kann
man die Laufzeit des Algorithmus nur mit einer
Wahrscheinlichkeitsbetrachung absch"atzen. Es sei $p$ der kleinste
Primfaktor von $n$. Die Laufzeit h"angt dann nur von der gr"o\3ten
Primpotenz von $p-1$ ab, sofern $p$, wie im allgemeinen zu erwarten
ist, als erster Faktor gefunden wird.

Betrachten wir $p-1$ als zuf"allig gew"ahlte nat"urliche Zahl, dann
hat der gr"o\3te Primfaktor im Mittel eine Gr"o\3enordung von
$(p-1)^{0.63}$. Falls also die gr"o\3te Primpotenz kleiner der
vorgegeben Schranke $B$ ist, so f"uhrt der Algorithmus zum Erfolg,
wenn die Primfaktoren weit genug auseinander liegen.

%%
%% ************** Faktorisierung mit elliptischen Kurven
%%
\subsection{Faktorisierung mit elliptischen Kurven}

Der Ansatz, mit elliptischen Kurven zu faktorisieren, geht auf
H.W. Lenstra zur"uck. Dieser Ansatz stellt in gewisser Hinsicht eine
Verallgemeinerung der $p-1$-Methode von Pollard dar. Der Aufwand beim
Faktorisieren mit elliptischen Kurven ist im wesentlichen auch von der
Gr"o\3enordnung des gesuchten Primfaktors abh"angig.


\subsubsection{Elliptische Kurven "uber K"orpern}

Es sei $\K$ ein K"orper. Im Hinblick auf die betrachteten Anwendungen
gen"ugt es K"orper $\K$ der Charakteristik $\mbox{char}(\K)\neq 2,3$
zu betrachten. Dann lassen sich alle elliptischen Kurven beschreiben
durch die Gleichungen der Form
\begin{equation} 
  y^2 = x^3 + Ax + B, 
\end{equation}
wobei $A,B\in\K$ mit $4A^3 + 27B^2 \neq 0$. Diese Nebenbedingung
sichert, da\3 das Polynom $x^3 + Ax + B$ keine doppelten Nullstellen
enth"alt.

F"ur $\K=\R$ lassen sich die elliptischen Kurven geometrisch
veranschaulichen. Interessant werden die
elliptischen Kurven $E(\R)$ dadurch, da\3 man auf ihnen in einem
abstrakten Sinne addieren kann. Zu je zwei Punkten $P_1$ und $P_2$ der
Kurve findet man durch eine geometrische Konstruktion stets einen
dritten, der ebenfalls auf der Kurve liegt und $P_1 + P_2$ genannt
wird. Die so definierte Addition folgt den "ublichen Rechenregeln: Sie
ist assozitativ, kommutativ, es existiert ein neutrales Element, und zu
jedem Punkt existiert ein Inverses. Eine elliptische Kurve bildet also
eine kommutative Gruppe bez"uglich dieser Addition.

Das Inverse $-P$ eines Kurvenpunktes $P$ ist dessen Spiegelbild an der
$x$-Achse. Die Summe zweier Punkte $P_1$ und $P_2$ findet man im
allgemeinen, indem man die Gerade durch die Punkte $P_1$ und $P_2$
zieht. Sie schneidet die Kurve in genau einem dritten Punkt
$Q$. Dessen Spiegelbild $-Q$ bez"uglich der $x$-Achse ist die gesuchte
Summe $P_1 + P_2$.

In dem Sonderfall $P_1=P_2$ tritt an die Stelle der Geraden durch
$P_1$ und $P_2$ die Tangente an die Kurve im Punkt $P_1=P_2$. Wenn
schlie\3lich $P_1=-P_2$ ist, trifft die Gerade durch $P_1$ und $P_2$
keinen weiteren Punkt der Kurve. Man stellt sich ersatzweise vor,
da\3 die Gerade die Kurve im Unendlichen schneidet.  Darum f"ugt man
noch einen weiteren Punkt hinzu, den man sich im Unendlichen denkt.
Man nennt ihn ${\cal O}$ und setzt $P_1-P_1={\cal O}$. Der Punkt ${\cal
O}$ ist das neutrale Element der additiven Gruppe $E(\R)$.

Die geometrische Defintion der Punktaddition kann in Formeln umgesetzt
werden. Wenn $P_1=(x_1,x_2)$ und $P_2=(x_2,y_2)$ Punkte der
elliptischen Kurve sind, dann gilt
\begin{equation} -P_1 = (x_1,-y_1) , 
\end{equation}
\begin{equation} P_1 + P_2 = 
                    \left\{ 
                       \begin{array}{cl}
                         {\cal O} & \mbox{ falls } P_1 = -P_2 \\
                         P_2      & \mbox{ falls } P_1 = {\cal O} \\
                         P_1      & \mbox{ falls } P_2 = {\cal O}
                       \end{array}
                    \right.
\end{equation}

In den anderen F"allen ist die Summe $P_3 = P_1 + P_2 = (x_3, y_3)$
definiert durch
\begin{equation}
\begin{array}{ccc}
  x_3 & = & \lambda^2 - x_1 - x_2                     \\
      &   &                                           \\
    y_3 & = & \lambda\left( x_1 - x_3\right) - y_1
\end{array}
\end{equation}
wobei 
$$ \lambda = \frac{y_1 - y_2}{x_1 - x_2} $$
falls $P_1\neq P_2$, und
$$ \lambda = \frac{3 x_1^2 + A}{2y_1}$$
falls $P_1 = P_2$.

Diese Formeln k"onnen nun eine Art Eigenleben gewinnen. Man darf sie
anwenden, ohne sich unter $x$ und $y$ Koordinaten von Punkten in der
Ebene vorstellen zu m"ussen. Es kann ein beliebiger
endlicher K"orper $\F$ verwendet werden. Geraden und Tangenten haben dann 
nicht mehr viel mit geometrischer Anschauung zu tun.

Betrachten wir nun n"aher die elliptischen Kurven "uber Primk"orpern
$\F_p$. Offensichtlich ist $E(\F_p)$ endlich. Man schreibt
"ublicherweise die Anzahl der Elemente in der Form
$$ \# E(\F_p) = p+1-t, $$ wobei $t$ als die Spur bezeichnet wird. F"ur
die Spur gilt die sogenannte Hasse-Schranke
$$ -2\sqrt{p} < t < 2\sqrt{p}. $$ 
Die Gruppe $E(\F_p)$ ist entweder eine zyklische Gruppe oder das
Produkt zweier zyklischer Gruppen.


%%
%% ************* "uber Ringen
%%

\subsubsection{Elliptische Kurven "uber Ringen}

Die Definition einer elliptischen Kurve "uber einem Restklassenring
ist ganz analog zur Definition der elliptischen Kurven "uber einem
K"orper.

\begin{definition}
  Eine elliptische Kurve "uber dem Restklassenring $\ring{n}$ ist definiert
  als die Menge der Paare $(x,y)\in\ring{n} \times \ring{n}$, die die
  Weierstrass-Gleichung
  $$ y^2 = x^3 + Ax + B $$
  mit $A$,$B\in\ring{n}$ und $4A^3 + 27B^2 \neq 0$ erf"ullen zusammen
  mit einem Punkt im Unendlichen ${\cal O}$. Sie wird mit
  $E(\ring{n})$ bezeichnet. Wir setzen dabei voraus, da\3 6 eine
  Einheit im Ring $\ring{n}$ ist.
\end{definition}

Auf dieser Punktmenge wird analog zur Addition auf elliptischen Kurven
"uber K"orpern eine Operation definiert. Eine direkte Verwendung der
oben definierten Additionsformeln ist nicht m"oglich, da die dort
auftretenden inversen Elemente nicht notwendigerweise existieren, da
$\ring{n}$ nicht nullteilerfrei ist, falls $n$ nicht prim ist.

Es sei zur Vereinfachung $n$ als quadratfrei vorausgesetzt, also $n=p_1
\cdot\ldots\cdot p_r$ f"ur paarweise verschiedene Primzahlen $p_i$
f"ur $i=1,\ldots,r$.  Dann existieren die Abbildungen
$$ h_i : E(\ring{n}) \longrightarrow E(\ring{p_i}), $$
die komponentenweise mittels der Projektionen
$$ \pi_i : \ring{n} \longrightarrow \ring{p_i} $$ 
f"ur $i = 1,\ldots,l$ definiert sind. Durch Zusammensetzen dieser
Abbildungen $h_i$ erh"alt man eine Abbilung
$$ h : E(\ring{n}) \longrightarrow E(\ring{p_1}) \times 
   \ldots \times E(\ring{p_l}), $$
die den gesamten Zielraum aussch"opft, bis auf alle Tupel, in denen das
neutrale Element einer der Kurven $E(\Z_{p_i})$ vorkommt und die nicht
das Nulltupel sind. Dabei
definieren wir die Abbildung $h$ so, da\3 das neutrale Element ${\cal
O}$ auf das Tupel bestehend aus den neutralen Elementen der Kurven
$E(\ring{p_1})$ bis $E(\ring{p_r})$ abgebildet wird. Der Chinesische
Restesatz liefert, da\3 allen Elementen im Wertebereich dieser
Funktion ein Urbild in eindeutiger Weise zugeordnet werden kann. F"ur
diese F"alle ist damit auch eine Operation auf $E(\ring{n})$ eindeutig
definiert, in dem man die zu verkn"upfenden Elemente zun"achst mittels
$h$ nach $E(\ring{p_1}) \times \ldots \times E(\ring{p_r})$ abbildet,
dort koordinatenweise addiert und gegebenfalls dem so erhaltenen Tupel
wieder sein Urbild unter $h$ zuordnet.

Die so definierte Operation kann auch durch Rechnen modulo $n$ direkt
in $\ring{n}$ ausgef"uhrt werden. Da diese Verkn"upfung, falls sie
definert ist, der Addition auf elliptischen Kurven "uber K"orpern
entspricht, bezeichnet man sie als {\em Pseudoadditon}. Die mit dieser
Operation versehene Kurve $E(\ring{n})$ bildet offensichtlich keine
Gruppe, da die Operation nicht total definiert ist. Es gilt f"ur jeden
Punkt $P\in E(\ring{n})$
$$ N_n \cdot P = {\cal O} \mbox{ mit } N_n := \kgv(\# E(\ring{p_1}),\ldots,\# E(\ring{p_r})). $$


\subsubsection{Algorithmus}

"Ahnlich wie bei den Primtests f"uhrt man beim Faktorisieren mit
elliptischen Kurven eine Berechnung aus, die gelingen w"urde, falls
$n$ eine Primzahl w"are.  Weil $n$ aber zusammengesetzt ist, kann die
Rechnung scheitern. Im Verlauf der Berechnung mu\3 man inverse Elemente
berechnen, indem man den gr"o\3ten gemeinsamen Teiler von $n$ mit
anderen Zahlen ausrechet. Die Rechnung kann nur dann fortgesetzt
werden, wenn der gr"o\3te gemeinsame Teiler gleich $1$ ist, also die
Zahlen teilerfremd sind. Falls der gr"o"ste gemeinsame Teiler dann auch
noch ungleich $n$ ist, hat man einen echten Teiler von $n$ gefunden.

Die Idee bei diesem Faktorisierungsalgorithmus ist es die Gruppe
$\einheit{n}$, die bei der $p-1$-Methode von Pollard verwendet wird,
durch die allgemeinere Gruppe der rationalen Punkte "uber den
elliptischen Kurven zu ersetzen. Auch hier ist der Algorithmus
erfolgreich, falls $\# E(\F_p)$ in kleine Primfaktoren zerf"allt. Der
wesentliche Vorteil liegt jedoch darin, da\3 zu einem vorgegeben $n$
verschiedene Gruppen mit verschiedener Ordungung aus dem Intervall
$$ \left[ p+1-2\sqrt{p},p+1+2\sqrt{p} \right] $$
zur Verf"ugung stehen. 

F"ur den Algorithmus kann man konkret wie folgt vorgehen. Zu einer
Schranke $b$ w"ahlt man "ahnlich wie bei der $p-1$-Methode von Pollard
einen Multiplikator $k:=b!$ oder einen "ahnlichen Multiplikator, der
nur aus Faktoren kleiner $b$ besteht. Dann w"ahlt man sich eine
elliptische Kurve $E(\ring{n})$ und einen Punkt $P$ auf der Kurve. In
der Praxis kann man dabei so vorgehen, da\3 man zun"achst nur den
Kurvenparameter $A$ vorgibt, dann einen Punkt $P$ zuf"allig w"ahlt und
dann den Kurvenparameter $B$ so bestimmt, da\3 der Punkt $P$ auf der
Kurve liegt. Dann berechnet man $k\cdot P$ auf der Kurve. Diese
Addition kann sehr effizient berechnet werden, da es nicht n"otig ist
$k$ mal $P$ zu addieren. Man kann "ahnlich wie bei "`square and
multiply"' vorgehen.  Die Addition f"uhrt man wie in der
definierten Pseudoaddition durch. Bei der Berechnung m"ussen Inverse
in $\ring{n}$ bestimmt werden, die in $\ring{n}$ nicht immer
existieren m"ussen. Ein Nullteiler hat einen nichttrivalen Teiler mit $n$. 
In diesem Fall liefert die Berechnung des gr"o\3ten gemeinsamen Teilers einen
Teiler von $n$.

Der Zufall spielt bei diesem Algorithmus wirklich mit, da es sehr viele
elliptische Kurven zur Auswahl gibt, an denen man die Berechnungen
durchf"uhren kann.  Lenstra konnte zeigen, da\3 es f"ur jedes
zusammengesetzte $n$ Kurven gibt, die einen Teiler liefern. Falls man
keinen nichttrivialen Teiler findet, kann man ein neues $k$ bestimmen
oder eine neue elliptische Kurve w"ahlen.

Der Aufwand dieses Faktorisierungsalgorithmus l"a\3t sich unter
gewissen bisher unbewiesenen heuristischen Annahmen wie folgt
absch"atzen. Es sei $p$ der kleinste Faktor der zu faktorisierenden
Zahl $n$. Es sei ferner $g$ eine beliebige nat"urliche Zahl. Dann
findet der Algorithmus mit Wahrscheinlichkeit $ \ge 1-e^{-g}$ den
Faktor $p$ und hat einen asymptotischen Aufwand von
$$ g \cdot \exp\left({\sqrt{(2+o(1))(\log p\log\log p)}} \left(\log n\right)^2\right) $$
f"ur $p\rightarrow \infty $.

Der Algorithmus eignet sich in hervorragender Weise zur
Parallelisierung. Man w"ahlt f"ur jeden zur Verf"ugung stehenden
Prozessor eine elliptische Kurve und rechnet auf allen Kurven solange
parallel, bis ein Faktor gefunden wurde. Dabei ist nur minimale
Vorbereitung und Nachbereitung auf jedem der Prozessoren notwendig.

%%
%% Faktorisierung mit quadratischen Kongruenzen
%%
\section{Faktorisierung mit quadratischen Kongruenzen}

Der Aufwand der Faktorisierungsalgorithmen mit Gruppen h"angt von der
Gr"o\3e der Primfaktoren der zu faktorisierenden Zahl $n$ ab.  F"ur
die Suche nach gr"o\3eren Teilern benutzt man andere Algorithmen, die
auf einer anderen Grundidee basieren. Die Laufzeit dieser Algorithmen
h"angt nur von $n$ ab.  Wir werden hier zwei Vertreter dieser
Algorithmenklasse kennenlernen.

Man will eine zusammengesetzte Zahl $n$ faktorisieren, indem man eine
quadratische Kongrunenz der Art
\begin{equation}
  x^2 \equiv y^2 \pmod{n} 
\end{equation}
aufstellt. Diese Kongruenz ist gleichbedeutend mit der Aussage, da\3
$n$ die Differenz $x^2 - y^2$ teilt. Nach der dritten binomischen
Formel ist
$$ x^2 - y^2 = (x-y)(x+y). $$ Wenn $n$ nicht selbst Teiler von $x-y$
oder $x+y$ ist, so mu\3 ein echter Teiler von $n$ in $x-y$ und ein
anderer echter Teiler in $x+y$ enthalten sein.  Durch die Berechnung
von $\ggt(x+y,n)$ kann dann ein nicht trivialer Faktor gefunden
werden.

Die Tatsache, da\3 sich "uber solche Kongruenzen ein Faktor finden
l"a\3t, ist schon l"anger bekannt. Sie l"a\3t sich praktisch aber nur
nutzen, wenn man geeignete Verfahren kennt, solche Kongruenzen zu
finden. Dies wurde f"ur nicht triviale F"alle erst durch die
Verwendung von leistungsf"ahigen Computern m"oglich.

Zun"achst werden wir den relativ einfachen Algorithmus von Dixon
betrachten, an dem sich eine ganze Reihe von grundlegenden Ideen
besonders gut erkl"aren lassen.  Er nimmt sofern eine Sonderstellung
ein, als da\3 f"ur diesen Algorithmus der subexponentielle Aufwand mit
mathematischer Strenge bewiesen werden kann, w"ahrend f"ur die
meisten anderen Varianten der Aufwand immer unter Verwendung
unbewiesener, lediglich auf heuristische Annahmen gest"utzte
Vermutungen abgesch"atzt wird.

%%
%% Grundlagen der Laufzeitanalyse
%%

\subsection{Grundlagen der Laufzeitanalyse}

Bevor wir uns den eigentlichen Algorithmen zuwenden, f"uhren wir
einige Definitionen ein und stellen einige S"atze vor, die uns die
Absch"atzung der Laufzeit der Faktorisierungsalgorithmen erleichtern
werden.

Wir benutzen im folgenden die Abk"urzung
\begin{equation} 
  L(n) := \exp\left(\sqrt{\log n\log\log n}\right). 
\end{equation}
Ferner benutzen wir die Abk"urzung $L^{\alpha}$ f"ur die Funktionenklasse
\begin{equation} L(n)^{\alpha + o(1)}. \end{equation}

\begin{lemma} F"ur diese Funktionenklasse gelten die folgenden Rechenregeln:
  \begin{itemize}
    \item $2 L^{\alpha} = L^{\alpha}$;
    \item $L^{\alpha} + L^{\beta} = L^{\max\{\alpha,\beta\}}$.
  \end{itemize}
\end{lemma}

\beweis 
"Ubung \hfill $\Box$

\begin{lemma} 
F"ur die Funktion $\pi(x)$, die die Anzahl der Primzahlen kleiner $x\in\R$ angibt, gilt
  \begin{equation}
    \pi(L^{\alpha}) = L^{\alpha} 
  \end{equation}
\end{lemma}

\beweis 
"Ubung \hfill $\Box$

Eine wichtige Formel, die bei der Laufzeitanalyse ben"otigt wird, ist
die Absch"atzung der Anzahl der nat"urlichen Zahlen, die nur
Primfaktoren kleiner einer vorgegebenen Schranke haben. Solche Zahlen
werden als {\em glatt} (engl. {\em smooth}) bzw. {\em $b$-glatt} bezeichnet,
wobei $b$ die Schranke f"ur die Gr"o\3e der Primfaktoren ist. Dieser
Begriff wurde von R.Rivest eingef"uhrt.

\begin{definition}
  Die Funktion $\glatt{x}{y}$ gibt die Anzahl der Elemente $n<x$ an,
  die glatt bez"uglich der Schranke $y$ sind:
  \begin{equation}
    \glatt{x}{y} := \#\{ n\in\N \mid n < x \mbox{ und } 
     (( p \mid n, p \mbox{ prim } ) \Rightarrow p < y )\}.
  \end{equation}
\end{definition} 

\begin{satz}
  F"ur eine beliebige Konstante $\epsilon$ sei $D:=\{(x,y) \mid x>10
  \mbox{ und } y \ge (\log x)^{1+\epsilon}\}$. Dann gilt f"ur alle
  Paare $(x,y)\in D$
  $$ \glatt{x}{y} = x \cdot u^{-u+o(u)} $$
  wobei $u:=\frac{\log x}{\log y}$.
\end{satz}

\begin{lemma} \label{spezialfall}
  F"ur den f"ur die Laufzeitanalyse wichtigen Spezialfall 
  $\glatt{n}{L^{\alpha}}$ gilt
  \begin{equation}
    \glatt{n}{L^{\alpha}} = n \cdot L^{-\frac{1}{2\alpha}}.
  \end{equation}
\end{lemma}

%%
%% Dixon Algorithmus
%%

\subsection{Dixons Algorithmus}

Der Algorithmus von Dixon ist f"ur praktische Zwecke nicht brauchbar,
aber verschiedene Grundideen k"onnen anhand dieses Algorithmus
besonders einfach erkl"art werden.  Au\3erdem handelt es sich um den
historisch ersten Algorithmus, f"ur den eine subexponentielle Laufzeit
exakt nachgewiesen werden konnte.

\subsubsection{Beschreibung des Algorithmus}

Es sei $n$ eine zusammengesetzte Zahl, die nicht durch Primzahlen
$p<L(n)$ teilbar ist. Zu $m\in\Z$ bezeichne $Q(m)$ den kleinsten
nicht-negativen Rest von $m^2 \pmod{n}$. Es sei $a$ eine Konstante $0
< a < 1$, die wir sp"ater optimal festlegen werden. Um eine
quadratische Kongruenz modulo $n$ aufzustellen, ben"otigen wir den
folgenden Teilalgorithmus:

%% ****** Teilalgorithmus
\begin{enumerate}  
  \item W"ahle $m\in\{1,\ldots,n\}$ zuf"allig aus und berechne $Q(m)$.

  \item Teste, ob $Q(m)$ nur Faktoren $p\le L^a$ besitzt. Falls dies
        der Fall ist, so faktorisiere $Q(m)$ vollst"andig und trage
        $m$, $Q(m)$ und das Tupel $(e_1,\ldots, e_r)$, wobei
        $Q(m)=\prod_{i=1}^{r} p_i^{e_i}$ und $p_1,\ldots,p_r$ die
        Primzahlen kleiner $L^a$ sind, in einer Tabelle ein.
\end{enumerate}

Dieser Teilalgorithmus wird so oft ausgef"uhrt, bis $\pi(L^a)+1$
Quadrate $Q(m)$, die $L^a$-glatt sind, gefunden wurden. Zu jedem
Eintrag haben wir einen Vektor
$$ v(m) \in\koerper{2}^{\pi(L^a)}, $$
dessen $i$-te Komponente eine 0 oder eine 1 ist, je nachdem, ob die
Primzahl $p_i$ in gerader oder in ungerader Potenz in $Q(m)$
auftritt.

Da wir $\pi(L^a)+1$ viele Vektoren der L"ange $\pi(L^a)$ haben,
existiert eine lineare Abh"angigkeit unter diesen Vektoren. Wir
rechnen "uber dem Skalark"orper $\koerper{2}$, das hei\3t uns
interessiert nur, ob der Exponent gerade oder ungerade ist. Nach
Umnummerierung k"onnen wir die lineare Abh"angigkeit als
$$ v(m_1)+v(m_2)+\ldots +v(m_k) = {\bf 0} .$$
ausdr"ucken. Es ist also das Produkt $Q(m_1) \cdot Q(m_2) \cdot \ldots
\cdot Q(m_k)$ ein Quadrat modulo $n$ und es
existiert folglich ein $x\in\Z$ mit
$$ Q(m_1) \cdot Q(m_2) \cdot \ldots \cdot Q(m_k) \equiv x^2 \pmod{n}. $$

Dieses $x$ kann sehr einfach mit Hilfe der im 2-ten Schritt
gespeicherten Tupel berechnet werden.  Ferner l"a\3t sich $y\equiv m_1
\cdot m_2 \cdot \ldots \cdot m_k \pmod{n}$ berechnen und es gilt
hierf"ur:
$$ y^2 \equiv m_1^2 \cdot m_2^2 \cdot \ldots \cdot m_k^2 \equiv Q(m_1) 
   \cdot Q(m_2) \cdot \ldots \cdot Q(m_k) \equiv x^2 \pmod{n} $$
Wir haben also eine quadratische Kongruenz modulo $n$
aufgestellt. Durch Berechnen von $\ggt(x+y,n)$ erhalten wir mit
Wahrscheinlichkeit gr"o\3er $1/2$ einen nichttrivialen Faktor von $n$.
Falls $x \neq \pm y \pmod{n}$ liefert die Berechnung von $\ggt(x+y,n)$
sicher einen nichttrivialen Faktor von n. Das folgende Lemma besagt
genauer:

\begin{lemma} 
  Es seien $x,y\in\ring{n}$ mit $x^2 \equiv y^2 \pmod{n}$
  gegeben. Dann folgt daraus mit Wahrscheinlichkeit $1-1/2^{t-1}$,
  da\3 $x \ne \pm y \pmod{n}$, wobei $t$ die Anzahl der Faktoren von
  $n$ ist.
\end{lemma}

\beweis 
"Ubung \hfill $\Box$

\begin{bem}
  Man sieht also, da\3 bei dem RSA-Verfahren die Verwendung von mehr
  als zwei Primzahlen keine Sicherheitsvorteile bringt.
\end{bem}

Um die Erfolgswahrscheinlichkeit bei diesem Algorithmus zu erh"ohen,
berechnet man anstatt $\pi(L^a)+1$ beispielsweise $\pi(L^a)+10$ solche
Eintr"age und erh"alt damit 10 verschiedene lineare Abh"angigkeiten,
von denen dann eine mit sehr hoher Wahrscheinlichkeit zu der gesuchten
Zerlegung von $n$ f"uhrt.

%%
%% Laufzeitanalyse
%%

\subsection{Laufzeitanalyse}

Wir f"uhren jetzt eine grobe Laufzeitanalyse durch, bei der auch der
Parameter $a$ optimal bestimmt wird.

In dem Teilalgorithmus m"ussen wir f"ur $\pi(L^a)=L^a$ viele
Primzahlen testen, ob sie $Q(m)$ teilen. Wir definieren $b:=b(n)$
indirekt dadurch, da\3 wir sagen, wir ben"otigen $L^b$ Durchl"aufe des
Teilalgorithmus, bis wir $\pi(L^a)+1$ viele geeignete Eintr"age
gefunden haben.  Dies ergibt einen Aufwand von $L^{a+b}$ Schritten,
bis wir gen"ugend viele gute $Q(m)$ gefunden haben.  Die
Gau\3elimination ben"otigt naiv implementiert $L^{3a}$ Schritte. Die
restlichen Berechnung der Werte $x$, $y$ und $\ggt(x+y,n)$ ben"otigt
nur $L^a$ viele Schritte.  Also kann die gesamte Laufzeit des
Algorithmus durch
$$ L^{\max\{ a+b,3a\}} $$
abgesch"atzt werden.

Wir m"ussen zur Bestimmung der Laufzeit die Funktion oder Konstante
$b$ ermitteln und $a$ optimal w"ahlen. Wir benutzen die heuristische
Annahme, da\3 bei einer zuf"alligen Wahl von $m$ auch das modulo $n$
reduzierte Quadrat $Q(m)$ zuf"allig verteilt ist. Damit betr"agt die
Wahrscheinlichkeit, da\3 $Q(m)$ in Faktoren $p < L^a$ zerf"allt,
$$ \frac{\glatt{n}{L^a}}{n}, $$
da $\glatt{n}{L^a}$ die Anzahl der $L^a$-glatten Zahlen kleiner $n$
angibt und $\frac{1}{n}$ unter der heuristischen Annahme die
Wahrscheinlichkeit jeder der m"oglichen Zahlen ist. Mit
Lemma~\ref{spezialfall} sehen wir, da\3
$$ \frac{\glatt{n}{L^a}}{n} = L^{-\frac{1}{2a}} $$
gilt. Um also $\pi(L^a) + 1 = L^a$ viele $L^a$-glatte $Q(m)$ 
zu finden, m"ussen wir
$$ L^a \cdot \left( L^{-\frac{1}{2a}} \right)^{-1} = L^{a+\frac{1}{2a}} $$
viele Werte f"ur $m$ untersuchen. Also gilt f"ur die Anzahl $b$ der
notwendigen Durchl"aufe des Teilalgorithmus
$$ b = a + \frac{1}{2a} , $$
und somit betr"agt der Aufwand
$$ L^{a+b} = L^{2a+\frac{1}{2a}}. $$
Die Funktion $2a + \frac{1}{2a}$ hat ihr Minimum bei $a=1/2$. Dies
ergibt eine Laufzeit f"ur den gesamten Algorithmus von
$$ L^{\max\{ a+b, 3a \}} = L^{\max\{ 2a + \frac{1}{2a}, 3a\}} = 
   L^{\max\{ 2, 1.5\}} = L^2, $$
das hei\3t es werden $L(n)^{2+o(1)}$ viele Schritte und ein
Speicheraufwand von $L(n)^{1+o(1)}$ f"ur die $L^a \times L^a$-Matrix
ben"otigt.

\begin{satz}
  Der Faktorisierungsalgorthmus von Dixon ben"otigt asymptotisch eine
  Laufzeit von $L(n)^{2+o(1)}$ vielen Schritten und einen
  Speicherplatzbedarf von $L(n)^{1+o(1)}$ bit.
\end{satz}

Unser Beweis dieses Satzes ist nicht mathmatisch exakt durchgef"uhrt
worden, da wir eine heuristische Annahme benutzt haben. Diese Annahme
betrifft die Wahrscheinlichkeit, mit der ein $Q(m)$ f"ur ein zuf"allig
gew"ahltes $m$ in Primfaktoren $p < L^a$ zerf"allt. Die $Q(m)$ sind nicht 
zuf"allig gleichverteilt, was wir angenommen haben.
Es l"a\3t sich jedoch zeigen, da\3 innerhalb einer Abweichung, 
die in $L(n)^{o(1)}$ abgefangen werden kann, dies auch f"ur die 
Quadrate $Q(m)$ zutrifft.

%%
%% Quadratischer Siebalgorithmus
%%

\subsection{Quadratischer Siebalgorithmus}

Ein weiterer wichtiger Ansatz, wie eine Kongruenz der Art
$$ x^2 \equiv y^2 \pmod{n} $$ 
aufgestellt werden kann, geht auf Pomerance zur"uck und verwendet ein
"ahnliches Verfahren, wie es beim Sieb des Eratosthenes eingesetzt
wird. Da mit Hilfe eines solchen Sieb-Verfahrens eine Quadratzahl
gefunden werden soll, spricht man von einem {\em quadratischen Sieb}.

In dem quadratischen Siebalgorithmus gehen wir von einem quadratischen
Polynom
\begin{equation} 
   f(x) = (x + \unten{\sqrt{n}})^2 - n \label{quadpolynom} 
\end{equation}
aus. F"ur dieses quadratische Polynom mit ganzzahligen Koeffizienten
gilt f"ur beliebige $x\in\Z$
$$ \left(x+\unten{\sqrt{n}}\right)^2 \equiv f(x) \pmod{n}. $$

In dem quadratischen Siebalgorithmus betrachten wir statt $Q(m)$ die
Werte $f(m)$, die in einer vorgegebenen Menge von $B$ Primzahlen
faktorisieren. Finden wir $m_1,\ldots,m_k \in\N$, so da\3
$$ f(m_1)\cdot \ldots \cdot f(m_k) $$
ein Quadrat ist, so erhalten wir mit
$$ x^2 := f(m_1)\cdot\ldots\cdot f(m_k) $$
und 
$$ y:=(m_1 + \unten{\sqrt{n}}) \cdot\ldots\cdot (m_k + \unten{\sqrt{n}}) $$
eine quadratische Kongruenz
$$ x^2 \equiv y^2 \pmod{n}, $$
mit der wir einen Faktor von $n$ bestimmen k"onnen. Die Zahl $x$ kann
wie beim Dixon-Algorithmus bestimmt werden, indem man die
gespeicherte Faktorisierung der $f(m_i)$ f"ur $i=1,\ldots,k$
benutzt.

Wir haben hierbei nur eine gewisse Wahrscheinlichkeit, da\3 diese so
gewonnene Gleichung wirklich zu einer L"osung f"uhrt, aber da wir mit
geringf"ugigem Aufwand mehrere solche Gleichungen aufstellen k"onnen,
l"a\3t sich die Wahrscheinlichkeit einen Faktor zu finden, leicht
erh"ohen. Damit haben wir also das Problem, $n$ zu faktorisieren, auf
das Problem, $k$ Zahlen $m_1$,\ldots,$m_k$ zu finden, f"ur die
$f(m_1)\cdot\ldots\cdot f(m_k)$ ein Quadrat modulo $n$ ist, reduziert.

Wie gro\3 sollte nun $B$ gew"ahlt werden? Eine geeignete Wahl f"ur $B$ 
ist sehr wichtig f"ur den Algorithmus. Wenn wir $B$ klein
w"ahlen, dann m"ussen wir weniger $m_i$ testen, f"ur die $f(m_i)$ in
kleine Primfaktoren zerf"allt, und das lineare Gleichungsystem kleiner.
Andererseits wird durch diese Wahl das Finden einer Linearkombination des 
Nullvektors schwieriger. Ein geeignetes $B$ stellt also einen Kompromis zwischen
diesen gegenl"aufigen Aspekten dar.

Heuristische "Uberlegungen legen nahe, da\3 $B$ in der Gr"o\3enordunung von
$$ \exp \left(\frac{1}{2}\sqrt{\log n \log\log n} \right) $$
gew"ahlt werden sollte.

Nehmen wir nun an, da\3 wir $B+10$ ganze Zahlen $m_1,\ldots,m_{B+10}$
finden k"onnen, f"ur die
$$ f(m_i) = \prod_{j=1}^{B} p_j^{\alpha_{ji}} $$
"uber der Faktorbasis $\{ p_1,\ldots p_B\}$ f"ur $1 \le i \le B+10$ 
faktorisieren. Die Exponenten $a_{ji}$ sind dabei nat"urliche Zahlen.

F"ur jedes $i=1,\ldots, B+10$ betrachten wir wie beim Dixon-Algorithmus
den Vektor
$$ v(m_i) = (\alpha_{1i},\ldots,\alpha_{B i}) \in\koerper{2}^{B} $$
Die Koeffizienten $\alpha_{ji}$ geben an, ob der Exponenten von $p_j$
in $f(m_i)$ gerade oder ungerade ist. Wir haben also $B+10$
Vektoren der L"ange $B$, unter denen wir eine Linearkombination
des Nullvektors suchen. Es existieren bis zu 10 verschiedene L"osungen, 
f"ur die gilt
$$ f(m_{i_1})\cdot\ldots\cdot f(m_{i_k}) = 
   \prod_{j=1}^B p_j^{\alpha_{j,i_1} + \ldots + \alpha_{j,i_k}} $$
und da
$$ \alpha_{j,i_1} + \ldots + \alpha_{j,i_k} \equiv 0 \pmod{2} $$
gilt, sind alle Exponenten gerade und somit ist
$$ f(m_{i_1}) \cdot\ldots\cdot f(m_{i_k}) $$
ein Quadrat. Damit liefert dieses Vorgehen also 10 potentielle
L"osungen, mit der eine quadratische Kongruenz aufgestellt werden
kann.

Die Methode in der jetztigen Form eignet sich noch nicht zur
Faktorisierung gro\3er Zahlen. Je kleiner die $f(m)$ sind, um so
gr"o\3er ist die Wahrscheinlichkeit, da\3 $f(m)$ "uber einer
vorgegebenen Faktorbasis $\{ p_1,\ldots,p_B\}$ zerf"allt. Daher
sollten wir m"oglichst kleine Werte f"ur $m$ nehmen, um geeignete
$f(m)$ zu finden.

Aber selbst f"ur $m$ aus dem Intervall $1\le m \le n^{\epsilon}$ mit
$\epsilon > 0$ und klein, liegen die Werte von $f(m)$ wegen
$$ f(m) = \left( m + \unten{\sqrt{n}}\right)^2 - n = 
   2m\unten{\sqrt{n}} + m^2 + \left( \unten{\sqrt{n}}^2 - n \right) 
   \approx 2m\unten{\sqrt{n}} + m^2 $$
in der Gr"o\3enordung $C\unten{\sqrt{n}}$ f"ur eine Konstante $C$, 
wenn $\epsilon < 1/4$ gilt.

Mit einem geeigneten Siebproze\3 k"onnen wir auf wesentlich
effektivere Weise geeignete Werte f"ur $m$ bestimmen. Wir
konzentrieren uns dabei nicht auf einzelne Werte f"ur $m$, sondern
betrachten gleich eine ganze Reihe von Werten.

Dazu ben"otigen wir jedoch noch einige Vor"uberlegungen. Zun"achst ist
klar, da\3 f"ur einen Primfaktor $p$ von $f(m_0)$ gilt, da\3 $p$ auch
Primfaktor von $ f(m_0 + k \cdot p) $ f"ur alle $k\in\Z$ ist. Um die
Menge der Werte f"ur $m$ zu bestimmen, f"ur die $p$ ein Faktor von
$f(m)$ ist, gen"ugt es folglich, die quadratischen Kongruenzen
$$ f(m) \equiv 0 \pmod{p} $$
f"ur $m\in\{0,\ldots, p-1\}$ zu l"osen.

Aus der speziellen Form des quadratischen Polynoms~\ref{quadpolynom}
sieht man, da\3 f"ur $p>2$, das nicht selbst Primfaktor von $n$ ist,
die Kongruenz
\begin{equation} 
   f(m) = (m + \unten{\sqrt{n}})^2 - n \equiv 0 \pmod{p} \label{siebkongruenz} 
\end{equation}
genau dann zwei verschiedene L"osungen besitzt, falls $n$ ein Quadrat
modulo $p$ ist, d.h. das Legendresymbol $\leg{n}{p}=1$. Falls $n$ kein
Quadrat modulo $p$ ist, d.h. das Legendresymbol $\leg{n}{p}=-1$,
existiert keine L"osung.  In diesem Fall ist $p$ kein Teiler von
$f(x)$ f"ur alle $x\in\N$. Wir bezeichnen daher mit
$p_1,\ldots,p_B$ nicht wie bei dem Dixon-Algorithmus die ersten $B$
Primzahlen, sondern die ersten $B$ Primzahlen $p$ mit $\leg{n}{p}=1$
inklusive 2. Diese Menge ist unsere {\em Faktorbasis}.

Die quadratische Siebprozedur funktioniert nun wie folgt. Zu jeder
Primzahl aus der Faktorbasis werden die beiden L"osungen 
\begin{eqnarray*}
  a_p & \equiv & n^{(p-1)/2} - \unten{\sqrt{n}} \pmod{p} \mbox{ und} \\
  b_p & \equiv & -n^{(p-1)/2} - \unten{\sqrt{n}} \pmod{p} 
\end{eqnarray*} 
Kongruenz~\ref{siebkongruenz} berechnet.  Wir betrachten nun nicht nur
einzelne zuf"allig gew"ahlte Werte f"ur $m$, sondern wir untersuchen
ein Intervall von $T$ vielen aufeinander folgenden $m$-Werten.

Zur Vorbereitung berechnen wir zun"acht f"ur jedes $m$ den Wert
$\unten{\log_2 f(m)}$, also einfach die Anzahl bits, die f"ur die
Bin"ardarstellung von $f(m)$ ben"otigt werden. Dieser Wert ist f"ur
viele aufeinanderfolgende $m$-Werte gleich und daher effizient zu
berechnen. Die $m$-Werte werden zusammen mit $\unten{\log_2 f(m)}$ in
einer Tabelle gespeichert.

Danach wird beim Sieben f"ur jede Primzahl $p$ aus der Faktorbasis der
Wert $\unten{\log_2 p}$ von dem Wert $\unten{\log_2 f(m)}$ abgezogen,
falls
$$ m \equiv a_p \pmod{p} \mbox{ oder } m \equiv b_p \pmod{p} $$
gilt.

F"ur ein $f(m)$, das h"ohere Potenzen einer Primzahl $p$ aus der
Faktorbasis enth"alt, werden noch zus"atzliche Tests ben"otigt. Das
hei\3t f"ur Primzahlen $p$ mit $p^k<C$ mu\3 f"ur $i<k$
$$ f(m) \equiv 0 \pmod{p^i} $$
getestet und gegebenfalls mehrfach $\unten{\log_2 p}$ abgezogen werden. 

Am Ende des Siebens wurde also von jedem Eintrag $\unten{\log_2 f(m)}$
f"ur jeden Primteiler $p$ von $f(m)$ der Wert $\unten{\log_2 p}$
entspechend oft abgezogen.  Falls also $f(m)$ "uber der Faktorbasis
zerf"allt, ist dieser Eintrag innerhalb der Rechenungenaugigkeit f"ur
die Rundung etwa gleich Null.

Mit diesem Sieb-Verfahren kann also aus einem Intervall von $T$ vielen
aufeinanderfolgenden Werten f"ur $m$ eine Liste von Werten $f(m)$
erstellt werden, die mit hoher Wahrscheinlichkeit "uber der
Faktorbasis zerfallen. Diese Werte werden dann noch einmal im
Einzelnen betrachtet, da w"ahrend des Siebverfahren die Faktoren nicht
protokolliert wurden. Aber jeder dieser Werte ergibt mit hoher 
Wahrscheinlichkeit einen Vektor "uber $\F_2$, so da"s
auf die oben skizzierte Art und Weise "uber die lineare
Abh"angigkeit ein Faktor von $n$ gefunden werden kann.

Eine auf heuristische Annahmen gest"utzte Laufzeitanalyse ergibt f"ur
diesen Algorithmus einen Aufwand von
$$ L(n)^{1+o(1)} = \exp\left( (1+o(1))\sqrt{\log n\log\log n} \right) .$$


   
   
      

      
      
    
    

    

    

